\scanpage{18}
\begin{center}
UNIVERSITE

DES SCIENCES ET TECHNIQUES

\uline{DU LANGUEDOC}
\end{center}

Service de la Scolarité

MA / GJ

Diffusion : Liste H

N° : 2152

Montpellier, le 28 avril 1983

\uline{AVIS DE SOUTENANCE D'UNE THESE DE DOCTORAT D'ETAT}

\uline{Nom et prénom de la candidate} : MIQUEL Geneviève épouse AMBLARD

\uline{Titre de la thèse} : ``Etude des mécanismes du transport transmembranaire des cations par un ionophore carboxylique, la grisorixine.

Etude expérimentale comparée sur des bicouches planes et des monocouches de lipides''.

\uline{Date, heure et lieu de soutenance prévus} : 13 MAI 1983 à 14h30 - Salle de Conférences

E.N.S.C. Montpellier - 8, rue de l'Ecole Normale

\uline{Nom du Directeur de thèse} : Monsieur Claude GAVACH

\uline{Désignation du laboratoire où la thèse est préparée} : LA.330. Laboratoire de Physico-chimie des Systèmes Polyphasés - C.N.R.S. Route de Mende - B.P. 5051 - 34033 MONTPELLIER CEDEX.

\uline{Résumé de la thèse} : Dans ce travail, on étudie les propriétés ionophores de la grisorixine dans deux systèmes physico-chimiques : monocouches mixtes ionophores-phospholipide à l'interface eau-air et bicouches planes de phospholipides (BLM) - solutions aqueuses.

Les méthodes expérimentales sont respectivement : manométrie superficielle et potentiel de surface - mesure des flux transmembranaire des ions ${}^{42}\mathrm{K}^{+}$ et des méthodes électrochimiques: mesures de résistance, potentiel de membrane et courant de court-circuit.

\uline{Durée approximative des recherches ayant abouti à la thèse} : 4 ans

\uline{Nom des Rapporteurs} :

M. BRUN Bernard, Professeur, Chimie Physique, U.S.T.L. Montpellier

M. GAVACH Claude, Maître de Recherches, Chimie Physique, C.N.R.S. Montpellier

M. PTAK Marius, Professeur, Chimie Physique, Faculté des Sciences d'Orléans.

\uline{Composition du jury proposé} :

M. BRUN Bernard, Chimie Physique, Professeur, U.S.T.L. Montpellier (Président)

M. GAVACH Claude, Chimie Physique, Maître de Recherches, C.N.R.S. Montpellier

M. HURWITZ H.D., Chimie Physique, Professeur, Faculté des Sciences, Bruxelles

(Membre extérieur choisi en raison de sa compétence scientifique)

M. PTAK Marius, Chimie Physique, Professeur, Faculté des Sciences- Orléans

M. TOCANNE Jean François, Biochimie, Maître de Recherches, C.N.R.S.- Toulouse

(Membre étranger à la spécialité de la candidate)

M. VERDIER Edgard, Chimie Physique, Professeur Associé, U.S.T.L.- Montpellier

M. LERNER Dan, Chimie Physique, Maître Assistant E.N.S.C.- Montpellier.

P. le Secrétaire Général,

l'Attachée d'Administration Universitaire,

\unsure{Gentil}

A. GENTIL
% archive foliation: 17 (inverted, upper right)
